Datasets for language models have rapidly expanded, culminating in the Common Crawl dataset\footnote{\url{https://commoncrawl.org/the-data/}} \cite{raffel2019t5} constituting nearly a trillion words.  This size of dataset is sufficient to train our largest models without ever updating on the same sequence twice. However, we have found that unfiltered or lightly filtered versions of Common Crawl tend to have lower quality than more curated datasets.  Therefore, we took 3 steps to improve the average quality of our datasets: (1) we downloaded and filtered a version of CommonCrawl based on similarity to a range of high-quality reference corpora, (2) we performed fuzzy deduplication at the document level, within and across datasets, to prevent redundancy and preserve the integrity of our held-out validation set as an accurate measure of overfitting, and (3) we also added known high-quality reference corpora to the training mix to augment CommonCrawl and increase its diversity.

Details of the first two points (processing of Common Crawl) are described in Appendix \ref{appendix:common_crawl_filtering}. For the third, we added several curated high-quality datasets, including an expanded version of the WebText dataset \cite{radford2019language}, collected by scraping links over a longer period of time, and first described in \cite{kaplan2020scaling}, two internet-based books corpora (Books1 and Books2) and English-language Wikipedia.


Table \ref{table:dataset} shows the final mixture of datasets that we used in training.  The CommonCrawl data was downloaded from 41 shards of monthly CommonCrawl covering 2016 to 2019, constituting 45TB of compressed plaintext before filtering and 570GB after filtering, roughly equivalent to 400 billion byte-pair-encoded tokens.    Note that during training, datasets are not sampled in proportion to their size, but rather datasets we view as higher-quality are sampled more frequently, such that CommonCrawl and Books2 datasets are sampled less than once during training, but the other datasets are sampled 2-3 times. This essentially accepts a small amount of overfitting in exchange for higher quality training data.

\begin{table}

    \begin{center}
        \begin{tabular}{lccc}
        \toprule
        Dataset & \shortstack{Quantity \\ (tokens)} & \shortstack{Weight in \\training mix} & \shortstack{Epochs elapsed when \\ training for 300B tokens} \\ 
        \midrule
        Common Crawl (filtered) & 410 billion & 60\% & 0.44\\ 
        WebText2 & 19 billion & 22\% &  2.9 \\
        Books1 & 12 billion & 8\% & 1.9 \\ 
        Books2 & 55 billion & 8\% & 0.43 \\ 
        Wikipedia & 3 billion & 3\% & 3.4 \\ 
        \bottomrule
        \end{tabular}
    \end{center}
    \caption{\textbf{Datasets used to train GPT-3}.  ``Weight in training mix'' refers to the fraction of examples during training that are drawn from a given dataset, which we intentionally do not make proportional to the size of the dataset.  As a result, when we train for 300 billion tokens, some datasets are seen up to 3.4 times during training while other datasets are seen less than once.}
    \label{table:dataset}
\end{table}

A major methodological concern with language models pretrained on a broad swath of internet data, particularly large models with the capacity to memorize vast amounts of content, is potential contamination of downstream tasks by having their test or development sets inadvertently seen during pre-training.  To reduce such contamination, we searched for and attempted to remove any overlaps with the development and test sets of all benchmarks studied in this paper.  Unfortunately, a bug in the filtering caused us to ignore some overlaps, and due to the cost of training it was not feasible to retrain the model. In Section \ref{section:measuring_and_preventing_memorization_of_benchmarks} we characterize the impact of the remaining overlaps, and in future work we will more aggressively remove data contamination. 